%% ============================================================
%%  LLM Introduction  –  v3: Reviewed & Revised Edition
%%  Changes from v2:
%%    - Technical corrections (n-gram count, GPT-4 tokens, RNN/PE claim,
%%      Q/K/V mechanics, PE range, multi-head motivation, RAG image,
%%      table title, reasoning claim)
%%    - Redundancy cuts (slides 21+22 merged, slides 37-41 removed,
%%      slide 26 removed, slides 48+49 merged, slide 60 removed)
%%    - New slides: Why LLMs Matter, Representation Evolution table,
%%      One Full Encoder Pass, SFT detail slide
%%    - Reorders: Sampling Strategies moved to §2; ChatGPT slide to §7
%%    - Fresher-accessibility rewrites throughout
%%  Sections:
%%    1.  Background & Motivation
%%    2.  Language Modeling Fundamentals
%%    3.  From NLP to Transformers
%%    4.  Transformer Architecture
%%    5.  Pre-training Paradigms (Encoder / Decoder / Enc-Dec)
%%    6.  LLM Training Pipeline
%%    7.  Inference & Use-Cases
%%    8.  LLM Landscape & Notable Models
%%    9.  Limitations, Ethics & Next Steps
%% ============================================================


%% ============================================================
%%  SECTION 1 – BACKGROUND & MOTIVATION
%% ============================================================

%%%%%%%%%%%%%%%%%%%%%%%%%%%%%%%%%%%%%%%%%%%%%%%%%%%%%%%%%%%%%%%%%%%%%%%%%%%%%%%%%%
\begin{frame}[fragile]\frametitle{}
\begin{center}
{\Large \textbf{Background \& Motivation}}\\[0.5em]
{\normalsize From Artificial Intelligence to Large Language Models}
\end{center}
\end{frame}

%%%%%%%%%%%%%%%%%%%%%%%%%%%%%%%%%%%%%%%%%%%%%%%%%%%%%%%%%%%
\begin{frame}[fragile]\frametitle{Why Large Language Models Matter to You}
\begin{itemize}
  \item LLMs are already embedded in tools you use: search, autocomplete,
    code assistants, customer support bots.
  \item Understanding how they work lets you use them better, evaluate
    their outputs critically, and build on top of them.
  % \item This course builds from first principles: by the end you will know
    % \emph{why} ChatGPT generates what it generates, not just \emph{that}
    % it does.
  \item No black-box thinking. Every component in the Transformer will be
    explained before it is used.
\end{itemize}
\end{frame}

%%%%%%%%%%%%%%%%%%%%%%%%%%%%%%%%%%%%%%%%%%%%%%%%%%%%%%%%%%%
\begin{frame}[fragile]\frametitle{Traditional Programming vs.\ Machine Learning}
\begin{center}
\includegraphics[width=0.6\linewidth,keepaspectratio]{tradml}
\end{center}
\begin{itemize}
\item \textbf{Traditional}: Human writes explicit rules $\rightarrow$ computer executes them.
\item \textbf{ML}: Human provides data + desired outputs $\rightarrow$ machine \emph{learns} the rules.
\item ML shines where hand-crafting rules is infeasible: vision, speech, language.
\end{itemize}
\end{frame}

%%%%%%%%%%%%%%%%%%%%%%%%%%%%%%%%%%%%%%%%%%%%%%%%%%%%%%%%%%%
\begin{frame}[fragile]\frametitle{AI $\supset$ ML $\supset$ Deep Learning $\supset$ LLMs}
\begin{center}
\includegraphics[width=0.82\linewidth,keepaspectratio]{ai1}
\end{center}
\begin{itemize}
\item \textbf{ML}: Statistical pattern learning from examples.
\item \textbf{Deep Learning}: Hierarchical feature learning via neural networks.
\item \textbf{LLMs}: Deep learning applied to language at planetary scale.
\end{itemize}
{\tiny (Ref: NVIDIA Blog – AI vs ML vs DL)}
\end{frame}

%%%%%%%%%%%%%%%%%%%%%%%%%%%%%%%%%%%%%%%%%%%%%%%%%%%%%%%%%%%
\begin{frame}[fragile]\frametitle{ML vs.\ Deep Learning: Feature Engineering}
\begin{center}
\includegraphics[width=\linewidth,keepaspectratio]{dlfeat1}
\end{center}
\begin{itemize}
\item Classic ML: human engineers features; model learns a mapping.
\item Deep Learning: model learns features \emph{and} mapping end-to-end from raw data.
\end{itemize}
{\tiny (Ref: XenonStack – ML vs DL)}
\end{frame}

%%%%%%%%%%%%%%%%%%%%%%%%%%%%%%%%%%%%%%%%%%%%%%%%%%%%%%%%%%%
\begin{frame}[fragile]\frametitle{The Rise of Large Language Models}

\begin{columns}
    \begin{column}[T]{0.5\linewidth}
		\begin{itemize}
		\item Scaling laws: performance improves predictably with more data, compute, and parameters.
		\item GPT-3 (175B params, 2020) showed remarkable few-shot capabilities.
		\item ChatGPT (2022) made LLMs accessible to hundreds of millions.
		\end{itemize}
		{\tiny (Ref: informationisbeautiful.net – Rise of Generative AI)}

    \end{column}
    \begin{column}[T]{0.5\linewidth}
		\begin{center}
		\includegraphics[width=\linewidth,keepaspectratio]{llm120}
		\end{center}
    \end{column}
  \end{columns}
  
  


\end{frame}

%%%%%%%%%%%%%%%%%%%%%%%%%%%%%%%%%%%%%%%%%%%%%%%%%%%%%%%%%%%
\begin{frame}[fragile]\frametitle{What is a Language Model?}
\begin{itemize}
\item While typing SMS, have you seen a next-word suggestion appear?
\item That is a \textbf{language model}: it assigns a probability to each possible next word given the words seen so far.
\item Think of it as: given everything said so far, what word is most likely to come next?
\item Modern LLMs extend this: they model \emph{long-range} dependencies across thousands of tokens using neural networks.
\item \textbf{Generative model}: keep sampling the next token until a complete response is formed.
\end{itemize}
\begin{center}
\includegraphics[width=0.55\linewidth,keepaspectratio]{chatgpt34}
\end{center}
\end{frame}

%%%%%%%%%%%%%%%%%%%%%%%%%%%%%%%%%%%%%%%%%%%%%%%%%%%%%%%%%%%
\begin{frame}[fragile]\frametitle{Generative vs.\ Discriminative Models}
\begin{columns}
  \begin{column}[T]{0.55\linewidth}
    \textbf{Discriminative} (e.g., logistic regression, BERT fine-tuned)
    \begin{itemize}
    \item Given input, predict label,  separates classes.
    \item Cannot generate new samples.
    \end{itemize}
    \vspace{0.4em}
    \textbf{Generative} (e.g., GPT, DALL-E)
    \begin{itemize}
    \item \textbf{Generative} models learn the full distribution of the data, not just ``which class does this belong to?'' but ``what does data in this domain look like?'' This lets them \emph{produce} new examples, not just classify existing ones.
    \item LLMs are generative: given a prompt, they sample a plausible continuation from everything they learned about how text is distributed.
    \end{itemize}
  \end{column}
  \begin{column}[T]{0.45\linewidth}
    \begin{center}
    \includegraphics[width=\linewidth,keepaspectratio]{llm64}
    \end{center}
    {\tiny (Ref: Avi Chawla – ML Community, Aug 2023)}
  \end{column}
\end{columns}
\end{frame}


%%%%%%%%%%%%%%%%%%%%%%%%%%%%%%%%%%%%%%%%%%%%%%%%%%%%%%%%%%%
\begin{frame}[fragile]\frametitle{Evolution of NLP Vectorisation}
\begin{center}
\includegraphics[width=\linewidth,keepaspectratio]{chatgpt30}
\end{center}
From simple frequency counts (bag-of-words, TF-IDF) $\rightarrow$ static embeddings (Word2Vec, GloVe) $\rightarrow$ contextualised representations (ELMo, BERT) $\rightarrow$ generative LLMs (GPT family).
{\tiny (Ref: Analytics Vidhya – Evolution of NLP)}
\end{frame}


%%%%%%%%%%%%%%%%%%%%%%%%%%%%%%%%%%%%%%%%%%%%%%%%%%%%%%%%%%%
\begin{frame}[fragile]\frametitle{What Problem Does Each Representation Solve?}
\begin{center}
\small
\renewcommand{\arraystretch}{1.3}
\begin{tabular}{|p{2cm}|p{2.5cm}|p{2.5cm}|}
\hline
\textbf{Method} & \textbf{Problem it solves} & \textbf{Problem it leaves} \\
\hline
Bag-of-Words & Words as countable symbols & Ignores word order \\
TF-IDF       & Weights rare, informative words & Still order-blind \\
Word2Vec     & Words as geometry (meaning $\approx$ distance) & One fixed vector per word \\
ELMo         & Context-sensitive vectors & Sequential; slow \\
Transformer  & Parallel + fully contextual & Needs huge data \& compute \\
\hline
\end{tabular}
\end{center}
\vspace{0.5em}
\begin{itemize}
  \item Each row fixes the row above's key limitation.
  \item The Transformer does not eliminate all problems, it trades old
    ones for new ones (scale, energy, hallucination).
\end{itemize}
\end{frame}


%% ============================================================
%%  SECTION 2 – LANGUAGE MODELING FUNDAMENTALS
%% ============================================================

% %%%%%%%%%%%%%%%%%%%%%%%%%%%%%%%%%%%%%%%%%%%%%%%%%%%%%%%%%%%%%%%%%%%%%%%%%%%%%%%%%%
% \begin{frame}[fragile]\frametitle{}
% \begin{center}
% {\Large \textbf{Language Modeling Prerequisite Concepts-Fundamentals}}\\[0.5em]
% {\normalsize How Next-Token Prediction Works}
% \end{center}
% \end{frame}


% %%%%%%%%%%%%%%%%%%%%%%%%%%%%%%%%%%%%%%%%%%%%%%%%%%%%%%%%%%%
% \begin{frame}[fragile]\frametitle{Statistical Approach: Count Next-Word Frequencies}
% \begin{itemize}
% \item Simple idea: for each word, count how often each other word follows it in the training corpus.
% \item Convert counts to probabilities via normalisation.
% \end{itemize}
% \begin{center}
% \includegraphics[width=0.7\linewidth,keepaspectratio]{llm2}

% \includegraphics[width=0.55\linewidth,keepaspectratio]{llm3}
% \end{center}
% {\tiny (Ref: What Is ChatGPT Doing\ldots – S.\ Wolfram)}
% \end{frame}

% %%%%%%%%%%%%%%%%%%%%%%%%%%%%%%%%%%%%%%%%%%%%%%%%%%%%%%%%%%%
% \begin{frame}[fragile]\frametitle{The n-gram Problem: Combinatorial Explosion}
% \begin{itemize}
% \item Using only the \textbf{current} word to predict the next: too little context, results are nonsensical.
% \item Using \textbf{2 words} (bi-gram): ${\sim}50{,}000^2 = 2.5$B possible pairs (${\sim}50{,}000$ is a typical sub-word vocabulary size, e.g.\ GPT-2 uses 50{,}257 tokens).
% \item Using \textbf{3 words} (tri-gram): ${\sim}125$ trillion combinations – storage and search nightmare.
% \item \textbf{Solution}: Learn a compact \emph{model} (parameterised neural network) that approximates the full distribution.
% \end{itemize}
% \begin{center}
% 
\includegraphics[width=0.75\linewidth,keepaspectratio]{llm11}
% \end{center}
% ChatGPT is precisely such a model – a \emph{large language model} (LLM) built to estimate next-token probabilities well.
% {\tiny (Ref: What Is ChatGPT Doing\ldots – S.\ Wolfram)}
% \end{frame}

%%%%%%%%%%%%%%%%%%%%%%%%%%%%%%%%%%%%%%%%%%%%%%%%%%%%%%%%%%%
\begin{frame}[fragile]\frametitle{What Is a Model? Parameters \& Knobs}
\begin{itemize}
\item A model replaces a lookup table with a compact mathematical function.
\item Example: a line $y = ax + b$ needs only 2 parameters to describe infinitely many $(x,y)$ pairs.
\item GPT-3 uses \textbf{175 billion parameters} – enough ``knobs'' to approximate the distribution of human language.
\item Training adjusts these parameters so the model assigns high probability to the actual next token in the training data.
\end{itemize}
\begin{center}
\includegraphics[width=0.7\linewidth,keepaspectratio]{llm12}
\end{center}
{\tiny (Ref: What Is ChatGPT Doing\ldots – S.\ Wolfram)}
\end{frame}

%%%%%%%%%%%%%%%%%%%%%%%%%%%%%%%%%%%%%%%%%%%%%%%%%%%%%%%%%%%
\begin{frame}[fragile]\frametitle{Tokens: The Unit LLMs Operate On}
\begin{itemize}
\item LLMs don't operate on words but on \textbf{tokens} (sub-word units).
\item Tokeniser (e.g.\ BPE, WordPiece) breaks text into pieces:
  \begin{center}
  \texttt{"playing"} $\rightarrow$ \texttt{["play", "\#\#ing"]}\\
  \texttt{"unbelievable"} $\rightarrow$ \texttt{["un", "\#\#believ", "\#\#able"]}
  \end{center}
\item Benefits: fixed vocabulary ($\sim$30K--50K tokens), handles rare/new words gracefully.
\item GPT-3: trained on 499B tokens; GPT-4: training data size undisclosed by OpenAI; estimates suggest well over 1T tokens.
\item In comparison: a human reading 12 hr/day for 80 years absorbs $\sim$5.3B words – a tiny fraction of GPT-3's training data alone.
\end{itemize}
\end{frame}


%%%%%%%%%%%%%%%%%%%%%%%%%%%%%%%%%%%%%%%%%%%%%%%%%%%%%%%%%%%
\begin{frame}[fragile]\frametitle{Next-Word Prediction}
\begin{itemize}
\item For the text \texttt{"The best thing about AI is its ability to"}, what comes next?
\item A language model produces a \textbf{ranked probability distribution} over all vocabulary words.
\item Sampling strategy matters:
  \begin{itemize}
  \item \textbf{Greedy (temp=0)}: always pick highest-prob word $\rightarrow$ flat, repetitive text.
  \item \textbf{Temperature 0.8}: occasionally pick lower-prob words $\rightarrow$ varied, creative output.
  \end{itemize}
\end{itemize}
\begin{center}
\includegraphics[width=0.6\linewidth,keepaspectratio]{llm1}
\end{center}
{\tiny (Ref: What Is ChatGPT Doing\ldots – S.\ Wolfram)}
\end{frame}

%%%%%%%%%%%%%%%%%%%%%%%%%%%%%%%%%%%%%%%%%%%%%%%%%%%%%%%%%%%
\begin{frame}[fragile]\frametitle{Sampling Strategies at Inference}
\begin{itemize}
\item \textbf{Greedy decoding}: always pick $\arg\max P(w)$ – fast, but repetitive and deterministic.
\item \textbf{Top-$k$ sampling}: sample from the $k$ most probable tokens (e.g.\ $k=50$).
\item \textbf{Top-$p$ (nucleus) sampling}: sample from the smallest set of tokens whose cumulative probability $\ge p$ (e.g.\ $p=0.9$).
\item \textbf{Temperature}: divide logits by $T$ before softmax.
  \begin{itemize}
  \item $T \to 0$: near-greedy; $T=1$: standard; $T > 1$: more random/creative.
  \end{itemize}
\item \textbf{Beam search}: maintain top-$b$ partial sequences; pick the one with highest joint probability.
\end{itemize}
\end{frame}

%% ============================================================
%%  SECTION 3 – FROM NLP TO TRANSFORMERS
%% ============================================================

%%%%%%%%%%%%%%%%%%%%%%%%%%%%%%%%%%%%%%%%%%%%%%%%%%%%%%%%%%%%%%%%%%%%%%%%%%%%%%%%%%
\begin{frame}[fragile]\frametitle{}
\begin{center}
{\Large \textbf{From NLP to Transformers}}\\[0.5em]
{\normalsize Embeddings, Seq2Seq, Attention}
\end{center}
\end{frame}

%%%%%%%%%%%%%%%%%%%%%%%%%%%%%%%%%%%%%%%%%%%%%%%%%%%%%%%%%%%
\begin{frame}[fragile]\frametitle{Word Vectors: Representing Words as Numbers}
\begin{columns}
  \begin{column}[T]{0.55\linewidth}
    \begin{itemize}
    \item Each word $\rightarrow$ a dense vector of $d$ real numbers (e.g.\ $d=300$).
    \item Geometrically similar words cluster together.
    \item Famous property: $\text{King} - \text{Man} + \text{Woman} \approx \text{Queen}$.
    \item Problem: static vectors – "bank" (river) and "bank" (financial) share one vector.
    \end{itemize}
  \end{column}
  \begin{column}[T]{0.45\linewidth}
    \begin{center}
    \includegraphics[width=0.65\linewidth,keepaspectratio]{bert7}
    \end{center}
  \end{column}
\end{columns}
{\tiny (Ref: CS224n – Christopher Manning)}
\end{frame}


%%%%%%%%%%%%%%%%%%%%%%%%%%%%%%%%%%%%%%%%%%%%%%%%%%%%%%%%%%%
\begin{frame}[fragile]\frametitle{Seq2Seq: Encoder-Decoder with RNN}
\begin{center}
\includegraphics[width=0.8\linewidth,keepaspectratio]{nlp14}
\end{center}
\begin{itemize}
\item \textbf{Encoder} reads the input sequence, compresses it into a fixed-size \emph{context vector}.
\item \textbf{Decoder} generates the output word-by-word, conditioned on that context vector.
\item \textbf{Problem}: single context vector is an information bottleneck for long sequences.
\end{itemize}
{\tiny (Ref: Deep Learning and NLP A-Z – K.\ Eremenko)}
\end{frame}

%%%%%%%%%%%%%%%%%%%%%%%%%%%%%%%%%%%%%%%%%%%%%%%%%%%%%%%%%%%
\begin{frame}[fragile]\frametitle{Attention Mechanism: Breaking the Bottleneck}
\begin{center}
\includegraphics[width=0.8\linewidth,keepaspectratio]{llm118}
\end{center}
\begin{itemize}
\item Instead of one context vector, the decoder \textbf{attends} to all encoder hidden states.
\item Each decoder step computes a weighted sum over encoder states,  the weight reflects how relevant each input word is right now.
\item Think of it like a spotlight: when generating the word ``bank'' in French, the model shines the spotlight back on the word ``bank'' in English.
\end{itemize}
{\tiny (Ref: CS447 Natural Language Processing – J.\ Hockenmaier)}
\end{frame}

%%%%%%%%%%%%%%%%%%%%%%%%%%%%%%%%%%%%%%%%%%%%%%%%%%%%%%%%%%%
\begin{frame}[fragile]\frametitle{Contextualized Embedding}
\begin{itemize}
\item \textbf{The problem with static embeddings}: Word2Vec assigns one
  fixed vector to each word regardless of context. The word
  \texttt{Apple} gets a single vector, an average of its ``fruit''
  and ``company'' meanings. That average is accurate for neither.
\item \textbf{Contextualised embeddings fix this}: instead of a fixed
  vector, each occurrence of \texttt{Apple} gets a \emph{different}
  vector shaped by the words around it. ``I love Apple Juice'' pulls
  the vector toward fruit; ``Apple posted record revenue'' pulls it
  toward technology company.
\item ``You shall know a word by the company it keeps.'' – J.\ R.\ Firth
\item Transformers achieve this through \textbf{self-attention}: each token's
  representation is updated by a weighted sum of all other tokens' \emph{Value}
  vectors. The weights come from the dot product of that token's \emph{Query}
  with every other token's \emph{Key}, so semantically related tokens
  contribute more.
\end{itemize}

\end{frame}


% %%%%%%%%%%%%%%%%%%%%%%%%%%%%%%%%%%%%%%%%%%%%%%%%%%%%%%%%%%%
% \begin{frame}[fragile]\frametitle{Contextualized Embedding}
% \begin{center}
% \includegraphics[width=0.72\linewidth,keepaspectratio]{llm256}
% \end{center}
% {\tiny (Ref: Deep Learning for Natural Language – Transformers – MIT OpenCourseWare)}
% \end{frame}


% %%%%%%%%%%%%%%%%%%%%%%%%%%%%%%%%%%%%%%%%%%%%%%%%%%%%%%%%%%%
% \begin{frame}[fragile]\frametitle{Self-Attention }

% \begin{center}
% \includegraphics[width=0.8\linewidth,keepaspectratio]{llm257}
% \end{center}

% {\tiny (Ref: Deep Learning for Natural Language – Transformers – MIT OpenCourseWare)}
% \end{frame}


%%%%%%%%%%%%%%%%%%%%%%%%%%%%%%%%%%%%%%%%%%%%%%%%%%%%%%%%%%%
\begin{frame}[fragile]\frametitle{Self-Attention: Every Token Attends to Every Token}
\begin{itemize}
\item \textbf{Cross-attention} (seq2seq): decoder attends to \emph{encoder} positions.
\item \textbf{Self-attention}: every position in a sequence attends to every \emph{other position in the same sequence}.
\item Enables the encoder to build rich, contextualised representations \emph{without} recurrence.
\item Enables full \textbf{parallelisation}: all tokens processed simultaneously during training,  dramatically faster than RNNs.
\end{itemize}
\vspace{0.5em}
Self-attention is the core innovation of the Transformer, replacing sequential RNN processing with parallel, content-based lookups.
\end{frame}



%% ============================================================
%%  SECTION 4 – TRANSFORMER ARCHITECTURE (DEEP DIVE)
%% ============================================================

%%%%%%%%%%%%%%%%%%%%%%%%%%%%%%%%%%%%%%%%%%%%%%%%%%%%%%%%%%%%%%%%%%%%%%%%%%%%%%%%%%
\begin{frame}[fragile]\frametitle{}
\begin{center}
{\Large \textbf{Transformer Architecture: Deep Dive}}\\[0.5em]
{\normalsize Step-by-step through each component}
\end{center}
\end{frame}


% %%%%%%%%%%%%%%%%%%%%%%%%%%%%%%%%%%%%%%%%%%%%%%%%%%%%%%%%%%%
% \begin{frame}[fragile]\frametitle{The Transformer: Overview}
% \begin{columns}
  % \begin{column}[T]{0.62\linewidth}
    % \begin{center}
    % \includegraphics[width=\linewidth,keepaspectratio]{promptengg81}
    % {\tiny (Ref: Prompt Engineering for AI Bootcamp, 2023)}
    % \end{center}
  % \end{column}
  % \begin{column}[T]{0.38\linewidth}
    % \begin{itemize}
    % \item Published: \emph{Attention Is All You Need} (Vaswani et al., 2017).
    % \item No recurrence $\Rightarrow$ fully parallel training.
    % \item Encoder stack (left) + Decoder stack (right).
    % \item Position captured by \textbf{Positional Encodings}.
    % \item Basis for BERT, GPT, T5, LLaMA, and all modern LLMs.
    % \end{itemize}
  % \end{column}
% \end{columns}
% \end{frame}



%%%%%%%%%%%%%%%%%%%%%%%%%%%%%%%%%%%%%%%%%%%%%%%%%%%%%%%%%%%
\begin{frame}[fragile]\frametitle{The Transformer: Overview}

    \begin{center}
    \includegraphics[width=\linewidth,keepaspectratio]{promptengg81}
    {\tiny (Ref: Prompt Engineering for AI Bootcamp, 2023)}
    \end{center}
	
 \emph{Attention Is All You Need} (Vaswani et al., 2017).

\end{frame}


% %%%%%%%%%%%%%%%%%%%%%%%%%%%%%%%%%%%%%%%%%%%%%%%%%%%%%%%%%%%
% \begin{frame}[fragile]\frametitle{Transformer Building Blocks}
% \begin{center}
% \includegraphics[width=0.6\linewidth,keepaspectratio]{llm112}
% \end{center}
% {\tiny (Ref: Generative AI with LLMs – Abhinav Kimothi)}
% \end{frame}

% %%%%%%%%%%%%%%%%%%%%%%%%%%%%%%%%%%%%%%%%%%%%%%%%%%%%%%%%%%%
% \begin{frame}[fragile]\frametitle{Transformer Building Blocks}

    % \begin{center}
    % \includegraphics[width=\linewidth,keepaspectratio]{promptengg81}
    % {\tiny (Ref: Prompt Engineering for AI Bootcamp, 2023)}
    % \end{center}
	
 % \emph{Attention Is All You Need} (Vaswani et al., 2017).

% \end{frame}

%%%%%%%%%%%%%%%%%%%%%%%%%%%%%%%%%%%%%%%%%%%%%%%%%%%%%%%%%%%
\begin{frame}[fragile]\frametitle{Step 1: Tokenisation \& Token Embedding}

\begin{columns}
    \begin{column}[T]{0.6\linewidth}
		\begin{itemize}
		\item \textbf{Tokenisation}: raw text split into sub-word tokens using BPE or WordPiece.
		\item \textbf{Token Embedding}: each token ID mapped to a learned dense vector via an embedding matrix.
		\item Embeddings encode \emph{semantic similarity}: words used in similar contexts get nearby vectors.
		\item They are \emph{learned} during pre-training, not hand-crafted.
		\end{itemize}
		
		{\tiny (Ref: What Is ChatGPT Doing\ldots – S.\ Wolfram)}

    \end{column}
    \begin{column}[T]{0.4\linewidth}
		\begin{center}
		\includegraphics[width=\linewidth,keepaspectratio]{llm13}
		\end{center}
    \end{column}
  \end{columns}
  
  
\end{frame}

%%%%%%%%%%%%%%%%%%%%%%%%%%%%%%%%%%%%%%%%%%%%%%%%%%%%%%%%%%%
\begin{frame}[fragile]\frametitle{Step 2: Positional Encoding, The Problem}
\begin{itemize}
\item The Transformer processes all tokens in parallel, its key
  advantage over RNNs. But parallelism means token order is no
  longer implicit.
\item Unlike RNNs, which encode order implicitly through sequential hidden-state
  updates, the Transformer's self-attention is \emph{permutation-equivariant}:
  shuffle the input tokens and the attention outputs shuffle identically.
  We must inject position information explicitly.
\item A raw integer index (0, 1, 2\ldots) is unsuitable: large values at long
  sequences would dominate the embedding and generalise poorly to lengths
  not seen in training.
\item Normalising indices to $[0,1]$ also fails: the step size
  between positions depends on sequence length, so the model
  cannot learn a consistent meaning of ``adjacent.''
\item \textbf{Requirements for a good positional signal}: bounded
  magnitude, unique per position, consistent across sequence lengths,
  and preservable by the model's arithmetic operations.
\end{itemize}


\end{frame}

%%%%%%%%%%%%%%%%%%%%%%%%%%%%%%%%%%%%%%%%%%%%%%%%%%%%%%%%%%%
\begin{frame}[fragile]\frametitle{Which function(s) to use?}
\begin{itemize}
\item If you use single sine function, ie $sine(i)$, then at multiple places you have same value (as it is periodical). So instead of single scalar value from one curve, overlay it with say cosine function on alternate tokens. But still, at some places even their addition can repeat. Then add again sine, cosine pairs with different frequencies. Till $L/2$ pairs. Like a Fourier transform. Summation is unique for each token.
\item The original Transformer uses sine/cosine waves at different frequencies,  each position gets a unique pattern/vector of numbers added to its embedding.
\item Think of it as a \textbf{unique address tag} for each position in the sentence.
\item  Note that even positions correspond to a sine function and odd positions correspond to cosine functions.
\item Added element-wise to the token embedding:
\[
  \mathbf{x}_i = \mathbf{e}_i + \text{PE}(i)
\]
\end{itemize}


\end{frame}

%%%%%%%%%%%%%%%%%%%%%%%%%%%%%%%%%%%%%%%%%%%%%%%%%%%%%%%%%%%
\begin{frame}[fragile]\frametitle{Positional Encoding Formula}

\begin{itemize}

\item Suppose you have an input sequence of length $L$ and require the position of the $k$th object within this sequence. The positional encoding is given by sine and cosine functions of varying frequencies::
\[
  \text{PE}(k, 2i)   = \sin\!\left(\frac{k}{n^{2i/d}}\right), \quad
  \text{PE}(k, 2i+1) = \cos\!\left(\frac{k}{n^{2i/d}}\right)
\]
\item $k$: position index of a token in the sequence, $0 \leq k \leq L-1$
\item $d$ : Dimension of the output embedding space
\item $i$: dimension pair index, $0 \leq i \leq d/2 - 1$, with a single value of $i$ mapping to both sine and cosine functions
\item $n$: User-defined scalar, set to 10000 by the author

\end{itemize}

\begin{center}
\includegraphics[width=0.5\linewidth,keepaspectratio]{llm253}
\end{center}


{\tiny (Ref: GA Gentle Introduction to Positional Encoding in Transformer Models, Part 1 –  Mehreen Saeed)}

\end{frame}



%%%%%%%%%%%%%%%%%%%%%%%%%%%%%%%%%%%%%%%%%%%%%%%%%%%%%%%%%%%
\begin{frame}[fragile]\frametitle{Positional Encoding Explanation}

\begin{center}
\includegraphics[width=0.7\linewidth,keepaspectratio]{llm253}
\end{center}

\textit{Plain English}: position $k$ in the sequence gets a vector of
$d$ numbers. Even-indexed dimensions use a sine wave and odd-indexed
dimensions use a cosine wave, each at a different frequency controlled
by $n^{2i/d}$. Because the frequencies differ, no two positions produce
the same vector, every position has a unique ``fingerprint'' that the
model can read.

{\tiny (Ref: GA Gentle Introduction to Positional Encoding in Transformer Models, Part 1 –  Mehreen Saeed)}

\end{frame}

%%%%%%%%%%%%%%%%%%%%%%%%%%%%%%%%%%%%%%%%%%%%%%%%%%%%%%%%%%%
\begin{frame}[fragile]\frametitle{Step 2: Positional Encoding}

Example: ``I am a robot'', with $n=100$ and $d=4$.


\begin{center}
\includegraphics[width=0.7\linewidth,keepaspectratio]{llm254}
\end{center}

{\tiny (Ref: GA Gentle Introduction to Positional Encoding in Transformer Models, Part 1 –  Mehreen Saeed)}

\end{frame}

%%%%%%%%%%%%%%%%%%%%%%%%%%%%%%%%%%%%%%%%%%%%%%%%%%%%%%%%%%%
\begin{frame}[fragile]\frametitle{Step 2: Positional Encoding}

Example: ``I am a robot'', with $n=100$ and $d=4$.


\begin{center}
\includegraphics[width=0.7\linewidth,keepaspectratio]{llm255}

\end{center}

In fact, the positional encoding matrix would be the same for any four-letter phrase with $n=100$ and $d=4$.
Range is always $[-1, 1]$ and sinusoid for each position is different, you have a unique way of encoding each position.

{\tiny (Ref: GA Gentle Introduction to Positional Encoding in Transformer Models, Part 1 –  Mehreen Saeed)}

\end{frame}

%%%%%%%%%%%%%%%%%%%%%%%%%%%%%%%%%%%%%%%%%%%%%%%%%%%%%%%%%%%
\begin{frame}[fragile]\frametitle{Positional Encoding: Why This Method?}
\begin{itemize}
\item \textbf{Why not just append the position number?}
  A raw integer like $42$ has no bounded magnitude – it would dominate the embedding
  and scale differently at test-time lengths never seen in training.
\item \textbf{Why sinusoids?} Three properties make them ideal:
  \begin{enumerate}
  \item \textbf{Bounded}: $\sin/\cos \in [-1, 1]$ – same scale as normalised embeddings.
  \item \textbf{Unique per position}: each position gets a distinct vector of
        $(\sin, \cos)$ values across frequencies.
  \item \textbf{Relative positions are learnable}: $\text{PE}(pos+k)$ is a
        \emph{linear function} of $\text{PE}(pos)$, so the model can learn
        to attend by relative offset, not just absolute position.
  \end{enumerate}
\item \textbf{Why add, not concatenate?}
  Addition keeps the vector dimension $d$ fixed.
  Concatenation would double $d$, doubling all subsequent weight matrices.
  Empirically, addition works just as well and is cheaper.
\item \textbf{Modern alternative – RoPE} (Rotary Position Embedding):
  encodes position directly into the $Q$–$K$ dot product so attention scores
  are naturally relative; used in LLaMA, Mistral, GPT-NeoX.
\end{itemize}
\end{frame}


%%%%%%%%%%%%%%%%%%%%%%%%%%%%%%%%%%%%%%%%%%%%%%%%%%%%%%%%%%%
\begin{frame}[fragile]\frametitle{Step 3: Self-Attention }

Make 3 copies.

\begin{center}
\includegraphics[width=0.7\linewidth,keepaspectratio]{llm259}
\end{center}

{\tiny (Ref: Deep Learning for Natural Language – Transformers – MIT OpenCourseWare)}
\end{frame}

%%%%%%%%%%%%%%%%%%%%%%%%%%%%%%%%%%%%%%%%%%%%%%%%%%%%%%%%%%%
\begin{frame}[fragile]\frametitle{Step 3: Self-Attention }

Got scores

\begin{center}
\includegraphics[width=0.8\linewidth,keepaspectratio]{llm260}
\end{center}

{\tiny (Ref: Deep Learning for Natural Language – Transformers – MIT OpenCourseWare)}
\end{frame}

%%%%%%%%%%%%%%%%%%%%%%%%%%%%%%%%%%%%%%%%%%%%%%%%%%%%%%%%%%%
\begin{frame}[fragile]\frametitle{Tunable Self-Attention }

While making 3 copies, multiply by tunable weight matrix for each copy, called $w_K, w_Q, w_V$

\begin{center}
\includegraphics[width=0.8\linewidth,keepaspectratio]{llm261}
\end{center}

{\tiny (Ref: Deep Learning for Natural Language – Transformers – MIT OpenCourseWare)}
\end{frame}



%%%%%%%%%%%%%%%%%%%%%%%%%%%%%%%%%%%%%%%%%%%%%%%%%%%%%%%%%%%
\begin{frame}[fragile]\frametitle{Tunable Self-Attention }

\begin{center}
\includegraphics[width=0.8\linewidth,keepaspectratio]{llm262}
\end{center}

{\tiny (Ref: Deep Learning for Natural Language – Transformers – MIT OpenCourseWare)}
\end{frame}

%%%%%%%%%%%%%%%%%%%%%%%%%%%%%%%%%%%%%%%%%%%%%%%%%%%%%%%%%%%
\begin{frame}[fragile]\frametitle{Self-Attention }

\begin{center}
\includegraphics[width=\linewidth,keepaspectratio]{llm263}
\end{center}

{\tiny (Ref: Self Attention in Transformers | Transformers in Deep Learning – Learn with Jay)}
\end{frame}

%%%%%%%%%%%%%%%%%%%%%%%%%%%%%%%%%%%%%%%%%%%%%%%%%%%%%%%%%%%
\begin{frame}[fragile]\frametitle{Step 3: Scaled Dot-Product Attention}
\begin{itemize}
\item $W^Q$, $W^K$, $W^V$ are \textbf{learned weight matrices},
  one set per attention head, updated during training. They project
  the same input token vector $\mathbf{X}$ into three different
  ``views'' of that token, each suited to a different role in the
  attention computation.
\item For each token, compute three projections:
\[
  \mathbf{Q} = \mathbf{X}W^Q, \quad \mathbf{K} = \mathbf{X}W^K, \quad \mathbf{V} = \mathbf{X}W^V
\]
\item Attention formula:
\[
  \text{Attention}(\mathbf{Q}, \mathbf{K}, \mathbf{V}) 
    = \text{softmax}\!\left(\frac{\mathbf{Q}\mathbf{K}^\top}{\sqrt{d_k}}\right)\mathbf{V}
\]
\item \textbf{In plain English}: compute how relevant each word is to every other word (dot product), normalise to get weights (softmax), then mix the information accordingly.
\item The $\sqrt{d_k}$ scaling prevents scores from getting too large and making the model overconfident.
\end{itemize}
\end{frame}

%%%%%%%%%%%%%%%%%%%%%%%%%%%%%%%%%%%%%%%%%%%%%%%%%%%%%%%%%%%
\begin{frame}[fragile]\frametitle{Q, K, V: Meaning and Intuition}
\textbf{Think of attention as a soft database lookup:}
\begin{center}
\small
\renewcommand{\arraystretch}{1.3}
\begin{tabular}{|l|l|p{5.5cm}|}
\hline
\textbf{Vector} & \textbf{Analogy} & \textbf{What it represents} \\
\hline
\textbf{Query} $Q_i$ & Search query & ``What information does token $i$ need?'' \\
\hline
\textbf{Key}   $K_j$ & Index tag    & ``What information does token $j$ offer?'' \\
\hline
\textbf{Value} $V_j$ & Actual data  & ``The content token $j$ will contribute if selected'' \\
\hline
\end{tabular}
\end{center}
\vspace{0.3em}
\textbf{Word-level example} for sentence \texttt{"The cat sat on the mat"}:
\begin{itemize}
\item Token \texttt{"sat"} issues a Query: \textit{``I'm a verb – who is my subject?''}
\item Token \texttt{"cat"} has a Key that scores highly against that Query.
\item The weighted Value of \texttt{"cat"} flows into the representation of \texttt{"sat"}.
\end{itemize}
\end{frame}

%%%%%%%%%%%%%%%%%%%%%%%%%%%%%%%%%%%%%%%%%%%%%%%%%%%%%%%%%%%
\begin{frame}[fragile]\frametitle{Q, K, V: Step-by-Step for One Word}
\textbf{Input sentence}: \texttt{"The bank can guarantee deposits"}\\[0.4em]
For the word \texttt{"bank"} (position $i = 2$):
\begin{enumerate}
\item \textbf{Project to Q, K, V}:
      $\mathbf{q}_2 = \mathbf{x}_2 W^Q,\quad
       \mathbf{k}_2 = \mathbf{x}_2 W^K,\quad
       \mathbf{v}_2 = \mathbf{x}_2 W^V$
\item \textbf{Score against every other token's Key}:
      $s_{2j} = \mathbf{q}_2 \cdot \mathbf{k}_j / \sqrt{d_k}$
      for $j = 1 \ldots T$
\item \textbf{Normalise scores}: $\alpha_{2j} = \text{softmax}(s_{2j})$
      \hfill {\small (sum to 1 across $j$)}
\item \textbf{Aggregate Values}:
      $\mathbf{z}_2 = \sum_j \alpha_{2j}\,\mathbf{v}_j$
\end{enumerate}
\vspace{0.3em}
After seeing ``deposits'' and ``guarantee'', attention weights for \texttt{"bank"} should
be high on those tokens $\Rightarrow$ \texttt{"bank"} gets a \emph{financial} representation.
In a sentence about a river, weights shift to ``river'', ``flow'' $\Rightarrow$
\texttt{"bank"} gets a \emph{geographical} representation.\\[0.3em]
This is how the Transformer achieves \textbf{contextualised embeddings}.
\end{frame}




%%%%%%%%%%%%%%%%%%%%%%%%%%%%%%%%%%%%%%%%%%%%%%%%%%%%%%%%%%%
\begin{frame}[fragile]\frametitle{Step 4: Multi-Head Attention}

\begin{center}
\includegraphics[width=0.8\linewidth,keepaspectratio]{llm258}
\end{center}

{\tiny (Ref: Deep Learning for Natural Language – Transformers – MIT OpenCourseWare)}
\end{frame}


%%%%%%%%%%%%%%%%%%%%%%%%%%%%%%%%%%%%%%%%%%%%%%%%%%%%%%%%%%%
\begin{frame}[fragile]\frametitle{Step 4: Multi-Head Attention}
\begin{itemize}
\item Not just a single word may mean different things (e.g. Apple or Light etc), but even a single full sentence can mean different things.
\item Example: ``She saw the man with the telescope''. It has two meanings.
\item Multiple self-attention layers capture different sentence/meaning patterns simultaneously, each head attending to a different type of linguistic relationship.
\end{itemize}
\end{frame}


%%%%%%%%%%%%%%%%%%%%%%%%%%%%%%%%%%%%%%%%%%%%%%%%%%%%%%%%%%%
\begin{frame}[fragile]\frametitle{Step 4: Multi-Head Attention}
\begin{itemize}
\item Run $h$ attention functions \textbf{in parallel}, each with independent $W^Q_l, W^K_l, W^V_l$ projections:
\[
  \text{MultiHead}(\mathbf{X}) = \text{Concat}(\text{head}_1, \ldots, \text{head}_h)\,W^O
\]
\item \textbf{Why multiple heads?} Each head specialises in a different type of relationship:
  \begin{itemize}
  \item Head 1: syntactic agreement (subject–verb)
  \item Head 2: co-reference (pronoun resolution)
  \item Head 3: positional proximity … and so on
  \end{itemize}
\item Think of it like multiple people reviewing the same sentence simultaneously, each looking for something different,  then combining their findings.
\item BERT-base uses $h=12$ heads; GPT-3 uses $h=96$.
\end{itemize}
\end{frame}


% %%%%%%%%%%%%%%%%%%%%%%%%%%%%%%%%%%%%%%%%%%%%%%%%%%%%%%%%%%%
% \begin{frame}[fragile]\frametitle{Multi-Head Attention: How It Works Mechanically}
% \begin{columns}
  % \begin{column}[T]{0.52\linewidth}
    % For each head $l = 1 \ldots h$:
    % \begin{enumerate}
    % \item Project $\mathbf{X}$ to lower-dimension Q, K, V:
          % $d_k = d_{\text{model}} / h$
          % (e.g.\ $768/12 = 64$ per head in BERT-base)
    % \item Compute scaled dot-product attention independently.
    % \item Each head looks at the \emph{same} input but through
          % \emph{different learned projection matrices}.
    % \end{enumerate}
    % After all heads:
    % \begin{itemize}
    % \item Concatenate $h$ outputs: shape $(T, h \cdot d_k) = (T, d_{\text{model}})$
    % \item Multiply by $W^O$ to mix heads: shape $(T, d_{\text{model}})$
    % \end{itemize}
  % \end{column}
  % \begin{column}[T]{0.48\linewidth}
    % \textbf{Why this design?}
    % \begin{itemize}
    % \item Each head has a smaller budget ($d_k \ll d_{\text{model}}$) – forced to
          % specialise on one type of signal.
    % \item Multiple heads = multiple ``questions'' asked of the same token simultaneously.
    % \item Final $W^O$ projection lets heads \emph{collaborate}: information from
          % different heads is mixed before being passed to the FFN.
    % \item Total computation is the same as single-head attention with $d_{\text{model}}$,
          % so no extra cost.
    % \end{itemize}
  % \end{column}
% \end{columns}
% \end{frame}

%%%%%%%%%%%%%%%%%%%%%%%%%%%%%%%%%%%%%%%%%%%%%%%%%%%%%%%%%%%
\begin{frame}[fragile]\frametitle{Step 5: Add \& Norm}
\textbf{Residual Connection + Layer Norm:}
\[
  \mathbf{h} = \text{LayerNorm}(\mathbf{x} + \text{MultiHead}(\mathbf{x}))
\]
\begin{itemize}
\item \textbf{Why gradients vanish}: during training, error signals
  flow backward through every layer via multiplication. In deep
  networks, repeatedly multiplying small numbers drives the signal
  toward zero, early layers receive almost no learning signal.
  The residual (skip) connection provides a direct addition path,
  so the gradient can travel straight to early layers without
  passing through every multiplication.
\item \textbf{Residual (skip) connection}: allows gradients to flow directly to earlier layers, preventing vanishing gradients.
\item \textbf{Layer Norm}: normalises across features for each token independently, stabilising training.
\end{itemize}

\end{frame}

%%%%%%%%%%%%%%%%%%%%%%%%%%%%%%%%%%%%%%%%%%%%%%%%%%%%%%%%%%%
\begin{frame}[fragile]\frametitle{Step 5:  Feed-Forward Network}

\textbf{Position-wise Feed-Forward Network (FFN):}
\begin{itemize}
\item Applied independently to \emph{each token} after attention.
\item Attention \emph{collects} information from other tokens; FFN \emph{processes} it,  adding non-linearity and depth to each token's representation.
\item Think of attention as gathering evidence, and FFN as reasoning about it.
\end{itemize}
\end{frame}


%%%%%%%%%%%%%%%%%%%%%%%%%%%%%%%%%%%%%%%%%%%%%%%%%%%%%%%%%%%
\begin{frame}[fragile]\frametitle{One Full Encoder Pass: Putting It Together}
For input sentence \texttt{``The bank can guarantee deposits''}:
\begin{enumerate}
  \item \textbf{Tokenise}: split into sub-word tokens via BPE.
  \item \textbf{Embed}: each token ID $\rightarrow$ dense vector $\mathbf{e}_i$.
  \item \textbf{Add positional encoding}: $\mathbf{x}_i = \mathbf{e}_i + \text{PE}(i)$.
  \item \textbf{Multi-head self-attention}: each token attends to all others;
    ``bank'' shifts toward ``guarantee''/``deposits''.
  \item \textbf{Add \& Norm}: residual connection stabilises gradients.
  \item \textbf{Feed-forward network}: each token representation processed
    independently.
  \item \textbf{Repeat} steps 4--6 for $N$ layers (e.g.\ $N=12$ for
    BERT-base).
  \item \textbf{Output}: contextualised vector per token, ready for a
    task head.
\end{enumerate}
\textit{Plain English: each layer refines what each token ``means'' given
its neighbours, stacking understanding layer by layer.}
\end{frame}


%% ============================================================
%%  SECTION 5 – PRE-TRAINING PARADIGMS
%% ============================================================

%%%%%%%%%%%%%%%%%%%%%%%%%%%%%%%%%%%%%%%%%%%%%%%%%%%%%%%%%%%%%%%%%%%%%%%%%%%%%%%%%%
\begin{frame}[fragile]\frametitle{}
\begin{center}
{\Large \textbf{Pre-training Paradigms}}\\[0.5em]
{\normalsize Encoder-only, Decoder-only, Encoder-Decoder, built differently from the start}
\end{center}
\end{frame}

%%%%%%%%%%%%%%%%%%%%%%%%%%%%%%%%%%%%%%%%%%%%%%%%%%%%%%%%%%%
\begin{frame}[fragile]\frametitle{Three Transformer Families: A Common Misconception}
\textbf{Are encoder-only models just encoder-decoder models with the decoder removed?}\\[0.6em]
\begin{center}
\large\color{red}$\times$ NO.\quad\color{black}They are architecturally different from day one.
\end{center}
\vspace{0.4em}
\begin{center}
\small
\renewcommand{\arraystretch}{1.3}
\begin{tabular}{|l|l|l|l|}
\hline
\textbf{Type} & \textbf{Analogy} & \textbf{What it does} & \textbf{Examples} \\
\hline
Encoder-only    & Highlighter   & Reads \& \emph{understands} text & BERT, RoBERTa \\
Decoder-only    & Story writer  & \emph{Continues} / generates text & GPT-4, LLaMA \\
Encoder-Decoder & Translator    & Reads one thing, \emph{writes} another & T5, BART \\
\hline
\end{tabular}
\end{center}
\vspace{0.4em}
\begin{itemize}
\item Parts are \textbf{never split} from a full model.
\item Each family is \textbf{designed, trained and deployed with only its own components}.
\item The architecture choice determines the training objective and best-fit tasks.
\end{itemize}
\end{frame}

% %%%%%%%%%%%%%%%%%%%%%%%%%%%%%%%%%%%%%%%%%%%%%%%%%%%%%%%%%%%
% \begin{frame}[fragile]\frametitle{Three Fundamental Transformer Variants}
% \begin{center}
% \includegraphics[width=\linewidth,keepaspectratio]{llm111}
% \end{center}
% The architecture choice determines what pre-training objective is used and what tasks the model naturally excels at.
% {\tiny (Ref: Generative AI with LLMs – Abhinav Kimothi)}
% \end{frame}

%%%%%%%%%%%%%%%%%%%%%%%%%%%%%%%%%%%%%%%%%%%%%%%%%%%%%%%%%%%
\begin{frame}[fragile]\frametitle{Paradigm 1: Encoder-Only (e.g., BERT)}
% \begin{columns}
  % \begin{column}[T]{0.5\linewidth}
    \textbf{Training objective}: Masked Language Modelling (MLM)
    \begin{enumerate}
    \item Take a sentence; randomly mask $\sim$15\% of tokens.
    \item Train the model to predict the original token at masked positions.
    \end{enumerate}
    \vspace{0.3em}
    \textbf{Example:}\\
    Input: \texttt{"The cat [MASK] on the [MASK]."}\\
    Target: \texttt{"sat"}, \texttt{"mat"}
    \vspace{0.3em}
    \begin{itemize}
    \item Sees \emph{both left and right} context,  truly bidirectional.
    \item No decoder exists at all,  \textbf{not capable of generating text}.
    \item Best for: \textbf{classification, NER, extractive QA, semantic search}.
    \item Models: \textbf{BERT}, RoBERTa, ELECTRA.
    \end{itemize}
  % \end{column}
  % \begin{column}[T]{0.5\linewidth}
    % \begin{center}
    % \includegraphics[width=\linewidth,keepaspectratio]{llm113}
    % \end{center}
    % {\tiny (Ref: Generative AI with LLMs: Abhinav Kimothi)}
  % \end{column}
% \end{columns}
\end{frame}

%%%%%%%%%%%%%%%%%%%%%%%%%%%%%%%%%%%%%%%%%%%%%%%%%%%%%%%%%%%
\begin{frame}[fragile]\frametitle{Paradigm 2: Decoder-Only (e.g., GPT)}
% \begin{columns}
  % \begin{column}[T]{0.5\linewidth}
    \textbf{Training objective}: Causal Language Modelling (CLM)
    \begin{itemize}
    \item Given tokens $w_1 \ldots w_{t-1}$, predict $w_t$.
    \item Uses \textbf{causal (masked) self-attention}: when predicting position $i$, only positions $1 \ldots i$ are visible,  future tokens are blocked.
    \item \textbf{Why block future?} At inference the future genuinely doesn't exist; training must mirror this to learn correctly.
    \end{itemize}
    \vspace{0.3em}
    No encoder. No cross-attention.
    \begin{itemize}
    \item Best for: \textbf{text generation, dialogue, code, reasoning, few-shot tasks}.
    \item Models: \textbf{GPT-4, LLaMA 3, Mistral, Falcon, Gemini}.
    \end{itemize}
  % \end{column}
  % \begin{column}[T]{0.5\linewidth}
    % \begin{center}
    % \includegraphics[width=\linewidth,keepaspectratio]{llm114}
    % \end{center}
    % {\tiny (Ref: Generative AI with LLMs: Abhinav Kimothi)}
  % \end{column}
% \end{columns}
\end{frame}

% %%%%%%%%%%%%%%%%%%%%%%%%%%%%%%%%%%%%%%%%%%%%%%%%%%%%%%%%%%%
% \begin{frame}[fragile]\frametitle{Causal Masking: Why Decoder Can Only See the Past}
% \textbf{The fundamental reason is training–inference consistency:}
% \begin{itemize}
% \item At \textbf{inference}, when generating token $i$, future token $i+1$ does not exist yet,  physically impossible to attend to it.
% \item At \textbf{training}, the full sequence is fed in one batch for efficiency.
% Without masking, position $i$ would trivially copy token $i+1$,  the model learns nothing.
% \item The causal mask enforces the \emph{same information constraint} in both phases.
% \end{itemize}
% \vspace{0.4em}
% \begin{center}
% \small
% \renewcommand{\arraystretch}{1.3}
% \begin{tabular}{|p{2cm}|p{2cm}|p{2cm}|}
% \hline
 % & \textbf{Encoder self-attn} & \textbf{Decoder self-attn} \\
% \hline
% Sees & All tokens (bidirectional) & Only past + current tokens \\
% Goal & Understand the full input as richly as possible & Predict only what could plausibly come next, given what has been written so far \\
% \hline
% \end{tabular}
% \end{center}
% \vspace{0.3em}
% Allowing the decoder to see future tokens would be \emph{data leakage},
% like giving a student the answer sheet during an exam.\\[0.3em]
% \textit{The encoder is a reader; the decoder is a writer who cannot peek at the next page.}
% \end{frame}


%%%%%%%%%%%%%%%%%%%%%%%%%%%%%%%%%%%%%%%%%%%%%%%%%%%%%%%%%%%
\begin{frame}[fragile]\frametitle{GPT Inference: What Happens After You Send a Prompt}
\textbf{Example prompt}: \texttt{"Translate to French: The sky is"}

\begin{enumerate}
\item \textbf{Tokenise}: prompt split into token IDs by the BPE tokeniser.
\item \textbf{Embed + positional encoding}: each ID $\rightarrow$ embedding vector; add positional tag.
\item \textbf{Forward pass through all $N$ decoder layers}:
      causal self-attention + FFN, repeated $N$ times (e.g.\ $N=96$ for GPT-3).
\item \textbf{Project to vocabulary}: final hidden state $\rightarrow$ logit vector of size $\sim$50K.
\item \textbf{Sample}: apply temperature + softmax; sample (or argmax) one token, e.g.\ \texttt{"bleu"}.
\item \textbf{Append \& repeat}: \texttt{"bleu"} is appended to the sequence; repeat until \texttt{<|endoftext|>}.
\end{enumerate}
\vspace{0.3em}
\textbf{Key-Value Cache}: Keys and Values for all previous tokens are cached.  Only the new token's Q, K, V are computed each step,  making autoregressive generation efficient.
\end{frame}


%%%%%%%%%%%%%%%%%%%%%%%%%%%%%%%%%%%%%%%%%%%%%%%%%%%%%%%%%%%
\begin{frame}[fragile]\frametitle{Paradigm 3: Encoder-Decoder (e.g., T5, BART)}
% \begin{columns}
  % \begin{column}[T]{0.5\linewidth}
    \textbf{Training objective}: Translation (T5) or denoising (BART)
    \begin{itemize}
    \item \textbf{Encoder} reads source text,  full bidirectional attention.
    \item \textbf{Decoder} generates the target text,  causal attention + \textbf{cross-attention} to encoder.
    \item Cross-attention: decoder's Q comes from decoder; K and V come from encoder,  bridging source to target.
    \end{itemize}
    \vspace{0.2em}
    \begin{itemize}
    \item Best for: \textbf{translation, abstractive summarisation, dialogue, code generation}.
    \item Models: \textbf{T5, BART}.
    \end{itemize}
  % \end{column}
  % \begin{column}[T]{0.5\linewidth}
    % \begin{center}
    % \includegraphics[width=\linewidth,keepaspectratio]{llm115}
    % \end{center}
    % {\tiny (Ref: Generative AI with LLMs: Abhinav Kimothi)}
  % \end{column}
% \end{columns}
\end{frame}

%%%%%%%%%%%%%%%%%%%%%%%%%%%%%%%%%%%%%%%%%%%%%%%%%%%%%%%%%%%
\begin{frame}[fragile]\frametitle{Pre-training Paradigms: Summary Table}
\begin{center}
\small
\renewcommand{\arraystretch}{1.3}
\begin{tabular}{|p{2cm}|p{2cm}|p{2cm}|p{2cm}|}
\hline
\textbf{Paradigm} & \textbf{Objective} & \textbf{Representative Models} & \textbf{Primary Use-Cases} \\
\hline
Encoder-Only & MLM: predict masked tokens & BERT, RoBERTa & Classification, NER, retrieval \\
\hline
Decoder-Only & CLM: predict next token & GPT-4, LLaMA, Mistral & Generation, dialogue, code \\
\hline
Encoder-Decoder & Span corruption \& reconstruct & T5, BART & Translation, summarisation \\
\hline
\end{tabular}
\end{center}

\vspace{0.4em}
\textbf{Trend}: large decoder-only models with instruction fine-tuning increasingly dominate even seq2seq tasks, as scale compensates for architecture.\\[0.3em]
\textbf{Key takeaway}: each family is built differently from scratch,  the architecture is the training objective is the use-case.
\end{frame}

%%%%%%%%%%%%%%%%%%%%%%%%%%%%%%%%%%%%%%%%%%%%%%%%%%%%%%%%%%%
\begin{frame}[fragile]\frametitle{Training Data: Scale and Diversity}
\begin{itemize}
\item GPT-3: trained on $\sim$499B tokens from WebText, BooksCorpus, Wikipedia, etc.
\item GPT-4: training data size undisclosed by OpenAI; estimates suggest well over 1T tokens, including code, scientific papers, and multilingual text.
\item LLaMA 3.1: 15.6T tokens.
\item \textbf{Data quality} matters as much as quantity: Common Crawl is huge but noisy; curated datasets (C4, The Pile) improve downstream performance.
\item Pre-training is the most expensive step:
  \begin{itemize}
  \item GPT-3: estimated \$4–12M in compute.
  \item LLaMA 3 405B: thousands of A100 GPU-days.
  \end{itemize}
\end{itemize}
\end{frame}


%% ============================================================
%%  SECTION 6: LLM TRAINING PIPELINE
%% ============================================================

% %%%%%%%%%%%%%%%%%%%%%%%%%%%%%%%%%%%%%%%%%%%%%%%%%%%%%%%%%%%%%%%%%%%%%%%%%%%%%%%%%%
% \begin{frame}[fragile]\frametitle{}
% \begin{center}
% {\Large \textbf{LLM Training Pipeline}}\\[0.5em]
% {\normalsize Pre-training $\rightarrow$ SFT $\rightarrow$ RLHF}
% \end{center}
% \end{frame}

% %%%%%%%%%%%%%%%%%%%%%%%%%%%%%%%%%%%%%%%%%%%%%%%%%%%%%%%%%%%
% \begin{frame}[fragile]\frametitle{Key Stages in LLM Development}
% \begin{enumerate}
% \item \textbf{Pre-training}: self-supervised on massive raw text corpus.\\
  % {\small High cost, long duration, high expertise required.}
% \item \textbf{Prompt Engineering}: craft effective prompts to steer the pre-trained model.\\
  % {\small No training cost; limited customisation.}
% \item \textbf{Supervised Fine-Tuning (SFT)}: train on small curated (prompt, response) pairs.\\
  % {\small Low–medium cost; teaches instruction-following.}
% \item \textbf{Reinforcement Learning from Human Feedback (RLHF)}: align model to human values.\\
  % {\small Medium–high cost; requires reward model.}
% \item \textbf{Deployment}: quantisation, distillation, caching optimisations.\\
  % {\small Reduces inference cost and latency.}
% \end{enumerate}
% \end{frame}

% %%%%%%%%%%%%%%%%%%%%%%%%%%%%%%%%%%%%%%%%%%%%%%%%%%%%%%%%%%%
% \begin{frame}[fragile]\frametitle{Supervised Fine-Tuning (SFT)}
% \begin{itemize}
% \item \textbf{Dataset}: human-curated (instruction, ideal response) pairs.
  % \begin{center}
  % \texttt{Instruction}: "Summarise the following text in one sentence."\\
  % \texttt{Response}: "The article discusses..."
  % \end{center}
% \item Train the model with supervised cross-entropy loss on these pairs.
% \item Transforms a \emph{next-token predictor} into an \emph{instruction-follower}.
% \item \textbf{Full fine-tuning}: update all parameters, expensive.
% \item \textbf{Parameter-Efficient Fine-Tuning (PEFT)}: update only a small subset (LoRA, prefix tuning, adapters), $<$1\% of parameters, comparable quality.
% \end{itemize}
% \end{frame}

% %%%%%%%%%%%%%%%%%%%%%%%%%%%%%%%%%%%%%%%%%%%%%%%%%%%%%%%%%%%
% \begin{frame}[fragile]\frametitle{LoRA: Low-Rank Adaptation}
% \begin{columns}
  % \begin{column}[T]{0.55\linewidth}
    % \begin{itemize}
    % \item Instead of updating the entire weight matrix $W$ (millions/billions of numbers), LoRA learns only a tiny \textbf{correction} $\Delta W \approx BA$ where $B$ and $A$ are small matrices.
    % \item Only $A$ and $B$ are trained; original weights $W$ are frozen.
    % \item Reduces trainable parameters by $100\times$ or more.
    % \item \textbf{QLoRA}: quantise the base model to 4-bit; apply LoRA on top.  Fine-tune a 65B model on a single consumer GPU.
    % \end{itemize}
  % \end{column}
  % \begin{column}[T]{0.45\linewidth}
    % \begin{itemize}
    % \item \textbf{Benefits}:
      % \begin{itemize}
      % \item Low memory footprint
      % \item Fast training
      % \item Multiple task-specific adapters for same base model
      % \item Cost-effective deployment
      % \end{itemize}
    % \end{itemize}
  % \end{column}
% \end{columns}
% \end{frame}

% %%%%%%%%%%%%%%%%%%%%%%%%%%%%%%%%%%%%%%%%%%%%%%%%%%%%%%%%%%%
% \begin{frame}[fragile]\frametitle{Reinforcement Learning from Human Feedback (RLHF)}
% \begin{center}
% \includegraphics[width=0.85\linewidth,keepaspectratio]{llm62}
% \end{center}
% Goal: align LLM behaviour with human values, \textbf{Helpful, Honest, Harmless (HHH)}.
% {\tiny (Ref: Katharina Koerner, LinkedIn, Aug 2023)}
% \end{frame}

% %%%%%%%%%%%%%%%%%%%%%%%%%%%%%%%%%%%%%%%%%%%%%%%%%%%%%%%%%%%
% \begin{frame}[fragile]\frametitle{RLHF: Step 1: Reward Model Training}

% \begin{columns}
    % \begin{column}[T]{0.6\linewidth}
		% \begin{itemize}
		% \item Sample multiple completions from the SFT model for each prompt.
		% \item Human annotators \textbf{rank} completions from best to worst.
		% \item Train a \textbf{Reward Model (RM)}: given (prompt, response), output a scalar score.
		% \item This score captures human preference without needing labels for every possible response.
		% \end{itemize}

		% {\tiny (Ref: Generative AI with LLM: Abhinav Kimothi)}

    % \end{column}
    % \begin{column}[T]{0.4\linewidth}
			% \begin{center}
			% \includegraphics[width=\linewidth,keepaspectratio]{rag19}
			% \end{center}
    % \end{column}
  % \end{columns}
  
% \end{frame}

% %%%%%%%%%%%%%%%%%%%%%%%%%%%%%%%%%%%%%%%%%%%%%%%%%%%%%%%%%%%
% \begin{frame}[fragile]\frametitle{RLHF: Step 2: PPO Optimisation}

% \begin{columns}
    % \begin{column}[T]{0.6\linewidth}
		% \begin{itemize}
		% \item Use \textbf{Proximal Policy Optimisation (PPO)} to update the LLM.
		% \item Think of the LLM as an agent learning to write better responses to get higher rewards from the Reward Model.
		% \item \textbf{KL penalty}: keeps updated model close to the SFT checkpoint, preventing reward hacking (the model gaming the score rather than genuinely improving).
		% \end{itemize}

		% {\tiny (Ref: Generative AI with LLM: Abhinav Kimothi)}

    % \end{column}
    % \begin{column}[T]{0.4\linewidth}
		% \begin{center}
		% \includegraphics[width=\linewidth,keepaspectratio]{rag22}
		% \end{center}
    % \end{column}
  % \end{columns}
  
% \end{frame}

% %%%%%%%%%%%%%%%%%%%%%%%%%%%%%%%%%%%%%%%%%%%%%%%%%%%%%%%%%%%
% \begin{frame}[fragile]\frametitle{Alternative: Direct Preference Optimisation (DPO)}
% \begin{itemize}
% \item PPO is complex and unstable. \textbf{DPO} (Rafailov et al., 2023) simplifies alignment.
% \item Directly optimise the LLM on preference pairs $(\text{chosen}, \text{rejected})$ without a separate reward model.
% \item \textbf{In plain terms}: given the same prompt, the model is shown a preferred response and a rejected response and trained to make the preferred one more likely,  directly, without the RL loop.
% \item Much simpler pipeline; increasingly preferred over RLHF for its stability and lower cost.
% \end{itemize}
% \end{frame}


%% ============================================================
%%  SECTION 7: USE-CASES
%% ============================================================

%%%%%%%%%%%%%%%%%%%%%%%%%%%%%%%%%%%%%%%%%%%%%%%%%%%%%%%%%%%%%%%%%%%%%%%%%%%%%%%%%%
\begin{frame}[fragile]\frametitle{}
\begin{center}
{\Large \textbf{Use Cases }}\\[0.5em]
{\normalsize How each architecture is deployed}
\end{center}
\end{frame}

%%%%%%%%%%%%%%%%%%%%%%%%%%%%%%%%%%%%%%%%%%%%%%%%%%%%%%%%%%%
\begin{frame}[fragile]\frametitle{LLM Real-World Use Cases}
\begin{center}
\includegraphics[width=0.88\linewidth,keepaspectratio]{llm121}
\end{center}
{\tiny (Ref: Applied LLMs Mastery 2024: Aishwarya Reganti)}
\end{frame}


%%%%%%%%%%%%%%%%%%%%%%%%%%%%%%%%%%%%%%%%%%%%%%%%%%%%%%%%%%%
\begin{frame}[fragile]\frametitle{Key Stages in LLM Development}
\begin{enumerate}
\item \textbf{Pre-training}: self-supervised on massive raw text corpus.\\
  {\small High cost, long duration, high expertise required.}
\item \textbf{Prompt Engineering}: craft effective prompts to steer the pre-trained model.\\
  {\small No training cost; limited customisation.}
\item \textbf{Supervised Fine-Tuning (SFT)}: train on small curated (prompt, response) pairs.\\
  {\small Low–medium cost; teaches instruction-following.}
\item \textbf{Reinforcement Learning from Human Feedback (RLHF)}: align model to human values.\\
  {\small Medium–high cost; requires reward model.}
\item \textbf{Deployment}: quantisation, distillation, caching optimisations.\\
  {\small Reduces inference cost and latency.}
\end{enumerate}
\end{frame}

% %%%%%%%%%%%%%%%%%%%%%%%%%%%%%%%%%%%%%%%%%%%%%%%%%%%%%%%%%%%
% \begin{frame}[fragile]\frametitle{What Is ChatGPT Really Doing?}
% \begin{itemize}
% \item ChatGPT is a \textbf{decoder-only} Transformer fine-tuned to follow instructions.
% \item At inference it runs the autoregressive loop:
  % ingest prompt tokens $\rightarrow$ predict next token $\rightarrow$ append $\rightarrow$ repeat.
% \item It has no symbolic knowledge, beliefs, or understanding in the human sense.
% \item It is ``merely'' pulling a coherent thread of text from the statistics of human language, yet the results are astoundingly human-like.
% \item Response quality is proportional to prompt quality.  Garbage in $\rightarrow$ garbage out.
% \end{itemize}
% {\tiny (Ref: What Is ChatGPT Doing\ldots, S.\ Wolfram; Understanding ChatGPT, G.\ Moore)}
% \end{frame}

%%%%%%%%%%%%%%%%%%%%%%%%%%%%%%%%%%%%%%%%%%%%%%%%%%%%%%%%%%%
\begin{frame}[fragile]\frametitle{Supervised Fine-Tuning (SFT): What Actually Happens}
\begin{itemize}
  \item The pre-trained model knows language but not how to
    \emph{follow instructions}. SFT teaches it that.
  \item Training data: human-written (prompt, ideal response) pairs
   , typically thousands to tens of thousands of examples.
  \item The model is trained to maximise the probability of the ideal
    response given the prompt, using the same cross-entropy loss as
    pre-training.
  \item Because the dataset is small relative to pre-training,
    \textbf{LoRA} (Low-Rank Adaptation) is standard: freeze most
    weights, train only small adapter matrices.
    \begin{itemize}
      \item Reduces trainable parameters from billions to millions.
      \item Fits on a single consumer GPU.
    \end{itemize}
  \item Output: an \emph{instruction-following} model (e.g.\
    \texttt{llama-3-8b-instruct}).
\end{itemize}
\end{frame}

%%%%%%%%%%%%%%%%%%%%%%%%%%%%%%%%%%%%%%%%%%%%%%%%%%%%%%%%%%%
\begin{frame}[fragile]\frametitle{Retrieval-Augmented Generation (RAG)}

    \begin{itemize}
    \item LLMs have static knowledge frozen at training time.
    \item \textbf{RAG} combines a retrieval system with an LLM generator.
    \item Steps:
      \begin{enumerate}
      \item Embed query with an encoder model.
      \item Search a \textbf{vector database} (Chroma, Pinecone, Weaviate) for relevant chunks.
      \item Inject retrieved context into the prompt.
      \item Decoder-LLM generates a grounded answer.
      \end{enumerate}
    \item Reduces hallucinations; enables up-to-date information.
    \end{itemize}

\end{frame}


% %%%%%%%%%%%%%%%%%%%%%%%%%%%%%%%%%%%%%%%%%%%%%%%%%%%%%%%%%%%
% \begin{frame}[fragile]\frametitle{Practical: Hugging Face APIs}
% \begin{lstlisting}
% import requests

% # Sentence completion (decoder-only: bloomz)
% API = "https://api-inference.huggingface.co/models/bigscience/bloomz"
% headers = {"Authorization": "Bearer <TOKEN>"}

% out = requests.post(API, headers=headers, json={
    % "inputs": "Sherlock Holmes is a",
    % "parameters": {"max_length": 60, "temperature": 0.8}
% }).json()
% print(out[0]["generated_text"])

% # Extractive QA  (encoder-only: roberta)
% QA_API = "https://api-inference.huggingface.co/models/deepset/roberta-base-squad2"
% out = requests.post(QA_API, headers=headers, json={
    % "inputs": {"question": "What is my job?",
               % "context": "My name is Yogesh and I am an AI Coach."}
% }).json()
% print(out["answer"])   # "AI Coach"
% \end{lstlisting}
% \end{frame}


%% ============================================================
%%  SECTION 8: LLM LANDSCAPE & NOTABLE MODELS
%% ============================================================

%%%%%%%%%%%%%%%%%%%%%%%%%%%%%%%%%%%%%%%%%%%%%%%%%%%%%%%%%%%%%%%%%%%%%%%%%%%%%%%%%%
\begin{frame}[fragile]\frametitle{}
\begin{center}
{\Large \textbf{LLM Landscape \& Notable Models}}\\[0.5em]
{\normalsize Closed, open, and specialised systems}
\end{center}
\end{frame}

%%%%%%%%%%%%%%%%%%%%%%%%%%%%%%%%%%%%%%%%%%%%%%%%%%%%%%%%%%%
\begin{frame}[fragile]\frametitle{Representative LLMs}
\begin{center}
\small
\renewcommand{\arraystretch}{1.3}
\begin{tabular}{|l|c|c|c|c|}
\hline
\textbf{Model} & \textbf{Params} & \textbf{Context} & \textbf{Multimodal} & \textbf{Open} \\
\hline
GPT-4o (OpenAI) & N/A & 128K & Yes & No \\
Claude 3.5 Sonnet (Anthropic) & N/A & 200K & Yes & No \\
Gemini 1.5 Pro (Google) & N/A & 2M & Yes & No \\
LLaMA 3.1 405B (Meta) & 405B & 128K & No & Yes \\
Mistral Large 2 & 123B & 128K & No & Yes \\
Mixtral 8x22B (MoE) & 141B & 64K & No & Yes \\
Phi-3 Medium (Microsoft) & 14B & 128K & No & Yes \\
\hline
\end{tabular}
\end{center}
Key trends: context windows growing; multimodality standard; open-source closing the gap.\\[0.3em]
{\footnotesize Model landscape evolves rapidly; verify current versions at \texttt{huggingface.co/models} before use.}
\end{frame}

%%%%%%%%%%%%%%%%%%%%%%%%%%%%%%%%%%%%%%%%%%%%%%%%%%%%%%%%%%%
\begin{frame}[fragile]\frametitle{Notable Closed-Source Ecosystems}
\begin{columns}
  \begin{column}[T]{0.5\linewidth}
    \textbf{ChatGPT / GPT-4o (OpenAI)}
    \begin{center}
    \includegraphics[width=\linewidth,keepaspectratio]{promptengg82}
    {\tiny (Ref: Prompt Engineering AI Bootcamp)}
    \end{center}
    API + Playground at \texttt{platform.openai.com}
  \end{column}
  \begin{column}[T]{0.5\linewidth}
    \textbf{Claude (Anthropic)}
    \begin{center}
    \includegraphics[width=\linewidth,keepaspectratio]{promptengg90}
    {\tiny (Ref: Prompt Engineering AI Bootcamp)}
    \end{center}
    Constitutional AI alignment; 200K context; API at \texttt{console.anthropic.com}
  \end{column}
\end{columns}
\end{frame}

%%%%%%%%%%%%%%%%%%%%%%%%%%%%%%%%%%%%%%%%%%%%%%%%%%%%%%%%%%%
\begin{frame}[fragile]\frametitle{Notable Open-Source Models}
\begin{itemize}
\item \textbf{LLaMA 3.1 (Meta)}: 8B, 70B, 405B variants.  Commercial licence available. Strong multilingual performance.
\item \textbf{Mistral / Mixtral}: Sliding window attention for long context; Mixture-of-Experts (MoE) architecture in Mixtral.
\item \textbf{Phi-3 (Microsoft)}: 14B model matching much larger closed models via high-quality training data.
\item \textbf{Falcon (TII UAE)}: 7B–40B, free commercial use, available on Hugging Face.
\item \textbf{Gemma (Google)}: small, efficient models optimised for on-device deployment.
\item Open-source advantage: privacy, customisation, no API costs, on-premises deployment.
\end{itemize}
\end{frame}

%%%%%%%%%%%%%%%%%%%%%%%%%%%%%%%%%%%%%%%%%%%%%%%%%%%%%%%%%%%
\begin{frame}[fragile]\frametitle{Beyond Standard Transformers: Emerging Architectures}
\begin{itemize}
\item \textbf{Mixture of Experts (MoE)}: Mixtral 8×22B routes each token to 2 of 8 expert FFN layers, reducing active parameters per forward pass while keeping total capacity high.
\item \textbf{Long-Context Techniques}: Sliding window attention (Mistral), RoPE position interpolation, Ring Attention for 1M+ tokens.
% \item \textbf{FlashAttention}: IO-aware exact attention implementation, $3\times$ faster, $10\times$ less memory, no approximation.
% \item \textbf{State Space Models (SSMs)}: Mamba, alternative to attention, linear-time inference, strong on long sequences.
\item \textbf{Multimodal LLMs}: GPT-4o, Gemini 1.5, process text, images, audio, video in a unified model.
\end{itemize}
\end{frame}


%% ============================================================
%%  SECTION 9: LIMITATIONS, ETHICS & NEXT STEPS
%% ============================================================

%%%%%%%%%%%%%%%%%%%%%%%%%%%%%%%%%%%%%%%%%%%%%%%%%%%%%%%%%%%%%%%%%%%%%%%%%%%%%%%%%%
\begin{frame}[fragile]\frametitle{}
\begin{center}
{\Large \textbf{Limitations, Ethics \& Next Steps}}\\[0.5em]
{\normalsize Responsible use of Large Language Models}
\end{center}
\end{frame}

%%%%%%%%%%%%%%%%%%%%%%%%%%%%%%%%%%%%%%%%%%%%%%%%%%%%%%%%%%%
\begin{frame}[fragile]\frametitle{LLM Limitations: What They Cannot Do}
\begin{itemize}
\item \textbf{No beliefs or knowledge}: LLMs are next-token predictors; they have no semantic understanding in the human sense.
\item \textbf{Hallucination}: will confidently generate plausible-sounding but factually wrong content.
\item \textbf{Autoregressive error accumulation}: a wrong early token propagates; no backtracking or self-correction.
\item \textbf{Context window limit}: conversations appear coherent until the log exceeds the context length; earlier parts are "forgotten."
\item \textbf{Static knowledge}: training has a cutoff date; model cannot access real-time information (unless augmented with RAG/tools).
\item \textbf{Uncertain reasoning ability}: standard LLMs produce
  plan-like outputs through pattern matching, not verified look-ahead.
  Newer reasoning models (OpenAI o1, DeepSeek-R1) extend
  chain-of-thought at inference time and show stronger multi-step
  problem solving, but whether this constitutes ``true reasoning''
  remains an active research question.
\end{itemize}
\end{frame}

%%%%%%%%%%%%%%%%%%%%%%%%%%%%%%%%%%%%%%%%%%%%%%%%%%%%%%%%%%%
\begin{frame}[fragile]\frametitle{LLM Challenges: Technical \& Safety}
\begin{center}
\includegraphics[width=\linewidth,keepaspectratio]{llm123}
\end{center}
{\tiny (Ref: Applied LLMs Mastery 2024: Aishwarya Reganti)}
\end{frame}

%%%%%%%%%%%%%%%%%%%%%%%%%%%%%%%%%%%%%%%%%%%%%%%%%%%%%%%%%%%
\begin{frame}[fragile]\frametitle{LLM Safety Concerns}
\begin{columns}
  \begin{column}[T]{0.5\linewidth}
    \textbf{Attack Vectors}
    \begin{itemize}
    \item Prompt injection \& jailbreaking
    \item Data poisoning during training
    \item Model extraction / IP theft
    \item Misinformation generation at scale
    \item Bias amplification from training data
    \item Privacy: memorisation of PII
    \end{itemize}
  \end{column}
  \begin{column}[T]{0.5\linewidth}
    \textbf{Mitigation Strategies}
    \begin{itemize}
    \item Red teaming \& adversarial testing
    \item Constitutional AI / RLAIF
    \item Output watermarking
    \item Input/output filters \& classifiers
    \item Differential privacy in training
    \item Responsible disclosure policies
    \end{itemize}
  \end{column}
\end{columns}
\end{frame}

%%%%%%%%%%%%%%%%%%%%%%%%%%%%%%%%%%%%%%%%%%%%%%%%%%%%%%%%%%%
\begin{frame}[fragile]\frametitle{Ethical Considerations}
\begin{itemize}
\item LLMs can generate racist, sexist, or harmful content, they reflect the biases in their training data.
\item Awareness and continuous monitoring of ethical implications is essential.
\item \textbf{Principles for responsible use}:
  \begin{itemize}
  \item Transparency: disclose AI involvement.
  \item Accountability: humans remain responsible for LLM outputs.
  \item Fairness: test for demographic disparities.
  \item Privacy: do not feed sensitive personal data into public APIs.
  \end{itemize}
\item \textbf{Always verify}: never blindly trust LLM output in high-stakes domains (medical, legal, financial, safety-critical).
\end{itemize}
\end{frame}

%%%%%%%%%%%%%%%%%%%%%%%%%%%%%%%%%%%%%%%%%%%%%%%%%%%%%%%%%%%
\begin{frame}[fragile]\frametitle{Emerging Trends (2025--2026)}
\begin{itemize}
\item \textbf{Small Language Models (SLMs)}: Phi-3, Gemma, high capability, on-device, low cost.
\item \textbf{Reasoning Models}: OpenAI o1/o3, DeepSeek-R1, extended chain-of-thought at inference time.
\item \textbf{Multimodal LLMs}: unified audio-image-video-text understanding.
\item \textbf{Agentic AI}: LLMs as orchestrators with tools, APIs, memory, and multi-step planning (LangChain Agents, AutoGPT, OpenAI Assistants).
\item \textbf{Continual learning}: models that update knowledge without full retraining.
\item \textbf{Efficient training}: DPO replacing PPO; LoRA and QLoRA as standard fine-tuning.
\end{itemize}
\end{frame}

%%%%%%%%%%%%%%%%%%%%%%%%%%%%%%%%%%%%%%%%%%%%%%%%%%%%%%%%%%%
\begin{frame}[fragile]\frametitle{Next Steps in Your LLM Journey}
\begin{columns}
  \begin{column}[T]{0.5\linewidth}
    \textbf{Essential Skills}
    \begin{itemize}
    \item Prompt engineering \& few-shot design
    \item Fine-tuning with LoRA / QLoRA
    \item RAG pipeline implementation
    \item Model evaluation (BLEU, ROUGE, BERTScore)
    \end{itemize}
    \textbf{Tools to Master}
    \begin{itemize}
    \item Hugging Face Transformers + PEFT
    \item LangChain / LlamaIndex
    \item Vector databases (Chroma, Pinecone)
    \item Ollama for local deployment
    \end{itemize}
  \end{column}
  \begin{column}[T]{0.5\linewidth}
    \textbf{Key Papers to Read}
    \begin{itemize}
    \item \emph{Attention Is All You Need} (2017)
    \item \emph{BERT}: Devlin et al.\ (2018)
    \item \emph{Language Models are Few-Shot Learners}: GPT-3 (2020)
    \item \emph{InstructGPT}: Ouyang et al.\ (2022)
    \item \emph{LoRA}: Hu et al.\ (2021)
    \item \emph{DPO}: Rafailov et al.\ (2023)
    \end{itemize}
    \textbf{Practice Projects}
    \begin{itemize}
    \item Build a chatbot with LangChain
    \item Fine-tune LLaMA-3 on a domain dataset
    \item Deploy a RAG pipeline with Chroma
    \end{itemize}
  \end{column}
\end{columns}
\end{frame}